\documentclass[11pt,a4paper]{article}

\usepackage[margin=2.2cm]{geometry}
\usepackage{amsmath,amssymb,booktabs,array,longtable,multirow}
\usepackage{graphicx}
\usepackage{xcolor}
\usepackage{hyperref}
\usepackage{enumitem}
\usepackage{microtype}
\usepackage{float}
\usepackage{caption}

\hypersetup{
  colorlinks=true,
  linkcolor=blue!55!black,
  urlcolor=blue!55!black,
  citecolor=blue!55!black
}
\setlist{nosep,leftmargin=*}
\newcommand{\AION}{\textsc{Aion}}
\newcommand{\HexCore}{\textsc{HexCore}}
\newcommand{\PhysX}{\textsc{PhysX}}
\newcommand{\MuJoCo}{\textsc{MuJoCo}}
\newcommand{\pass}{\textcolor{green!45!black}{\textbf{PASS}}}
\newcommand{\reject}{\textcolor{red!65!black}{\textbf{REJECT}}}
\newcommand{\bounded}{\textcolor{orange!70!black}{\textbf{BOUNDED}}}

\title{From Motor Calibration to Predictive Franka Manipulation\\
\large AION Precision-Control Apprenticeship, Articulated Transfer, and NVIDIA PhysX Evaluation}
\author{AION / Tessaris Research Programme}
\date{4 August 2026}

\begin{document}
\maketitle

\begin{abstract}
This document records the complete development sequence by which \AION{}
progressed from generic stop-at-target calibration to articulated cube lifting
under local \MuJoCo{} physics and then to its first teacher-removed cube lift
under NVIDIA Isaac/\PhysX{}.  The programme did not treat vision accuracy,
joint imitation, or offline action error as manipulation success.  Each
candidate was required to commit actions before execution, receive only RGB
and bounded robot proprioception at runtime, and face independently executed
physical consequences.  A cross-body precision method reduced final local
positioning error by $1{,}620.37\times$, a local articulated controller lifted
4/4 sealed cubes against 0/4 for cold control, and the subsequent cloud
programme falsified two inadequate representations before a predictive
position--orientation controller completed one fresh \PhysX{} lift.  That
controller remained unpromoted because it scored only 1/12, below the fixed
greater-than-2/12 reliability gate.  The result is therefore a documented
advance in earned physical competence and failure localization, not a claim of
general robotic manipulation, NVIDIA world-model mastery, AGA, or AGI.
\end{abstract}

\section{Research Question and Operating Principle}

The original failure was often described as a cube-lifting problem, but the
deeper capability deficit was general movement judgment:

\begin{quote}
Given a target, an unfamiliar body, delayed and uncertain actuation, and only
partial visual feedback, can \AION{} predict the effect of a movement, brake
before overshoot, observe the resulting error, and issue progressively smaller
corrections until a physical objective is achieved?
\end{quote}

This is the shared structure behind reaching, grasping, walking, steering,
parking, and other embodied tasks.  Perception alone is insufficient.  The
controller must translate distance, direction, speed, orientation, timing, and
uncertainty into bounded actions whose real consequences are repeatedly
measured.

The retained precision loop is
\begin{equation}
\text{intervene}
\rightarrow \text{identify dynamics}
\rightarrow \text{predict}
\rightarrow \text{approach}
\rightarrow \text{brake}
\rightarrow \text{observe}
\rightarrow \text{micro-correct}.
\label{eq:precision-loop}
\end{equation}

The physical programme separates inexpensive practice from independent
examination:
\begin{equation}
\underbrace{\text{large-scale local \MuJoCo{} practice}}_{\text{cheap, varied, inspectable}}
\rightarrow
\underbrace{\text{frozen Isaac/\PhysX{} examination}}_{\text{paid, simulator-disjoint authority}}
\rightarrow
\text{diagnosis and revised local learning}.
\label{eq:practice-exam}
\end{equation}

NVIDIA execution is therefore not used for routine trial and error.  It is the
independent physics authority that detects whether a method learned genuine
control principles or merely exploited the local simulator.

\section{Authority Separation and Runtime Contract}

Every stage distinguishes proposal inputs from evaluator authority.

\begin{table}[H]
\centering
\caption{Runtime and evaluation authority boundary.}
\begin{tabular}{>{\raggedright\arraybackslash}p{0.27\linewidth}
                >{\raggedright\arraybackslash}p{0.30\linewidth}
                >{\raggedright\arraybackslash}p{0.33\linewidth}}
\toprule
Layer & Permitted & Withheld or evaluator-owned \\
\midrule
Local precision lanes & position/velocity observations, bounded actions, time
& hidden delay, actuator matrix, damping, sealed target dynamics \\
Local articulated arm & RGB, bounded joints, bounded gripper state, time
& cube pose, contact flags, physical lift score, teacher actions \\
Isaac/\PhysX{} & RGB, 18-dimensional bounded robot proprioception, episode step
& cube coordinates, simulator contacts, reward internals, teacher actions \\
Training only & quarantined geometry and successful teacher trajectories
& no authority to appear in the deployed runtime \\
Promotion & immutable receipts, cold/champion controls, fixed gates
& no score self-awarded by the proposal policy \\
\bottomrule
\end{tabular}
\end{table}

Every physical action is hashed and committed before execution.  Non-finite,
out-of-range, incorrectly dimensioned, or authority-expanding programs are
rejected.  A development success can authorize a later sealed exam but cannot
promote itself.

\section{Stage I: Cross-Body Precision-Control Apprenticeship}

\subsection{Identified forward dynamics}

For a previously unseen control body, \AION{} performs bounded interventions
to identify action delay $d$, actuation matrix $A$, and damping matrix $D$:
\begin{align}
\dot{v}_{t} &= A u_{t-d} - Dv_t, \\
q_{t+1} &= q_t + v_t\Delta t
          + \tfrac{1}{2}\dot{v}_{t}\Delta t^2.
\label{eq:forward-model}
\end{align}

The inverse controller uses the identified model to approach the target, while
a braking model reduces command magnitude near the stopping boundary.
Uncertainty reduces control authority.  Observations are incorporated after
every action rather than after a long open-loop trajectory.

The existing AION prediction engine remains responsible for executive-level
feasibility and risk forecasts.  The precision forward model is a faster,
numerical component for immediate actuator consequences.  These layers are
complementary rather than duplicates.

\subsection{Sealed result}

Six source-disjoint worlds varied control dimension, delay, strength, damping,
target, and post-action disturbance.  The sealed cohort ranged from one to
five dimensions.

\begin{table}[H]
\centering
\caption{Cross-body precision-control results.}
\begin{tabular}{lr}
\toprule
Measure & Result \\
\midrule
Learned controller success & 6/6 \\
Matched uncalibrated control & 0/6 \\
Mean learned final error & $0.000029995$ \\
Mean cold final error & $0.048603$ \\
Final-error improvement & $1{,}620.37\times$ \\
Hidden action-delay identification & 100\% \\
Disturbance recovery & 6/6 \\
Control dimensions & 1--5 \\
Pre-action commitments & 1,560 \\
Malicious action programs rejected & 3/3 \\
Near-uncontrollable actuator & explicit abstention \\
Memorized trajectories retained & 0 \\
Body-specific parameters retained & 0 \\
Restart reconstruction & pass \\
\bottomrule
\end{tabular}
\end{table}

\HexCore{} promoted
\texttt{procedure\_cross\_body\_precision\_control\_apprenticeship\_v1}.
The retained asset is the method
\texttt{identify\_predict\_brake\_micro\_correct}, not the sealed trajectories
or individual body parameters.

\section{Stage II: Local Articulated Precision and Cube Lifting}

\subsection{Composed control stack}

The precision method was composed with previously retained visual topology and
contact/occlusion skills:
\begin{equation}
\begin{split}
\text{RGB landmarks}
&\rightarrow \text{articulated topology}
\rightarrow \text{pixel-space inverse kinematics} \\
&\rightarrow \text{motor identification and relevance routing}
\rightarrow \text{approach} \\
&\rightarrow \text{opposing-finger contact}
\rightarrow \text{lift}
\rightarrow \text{physical outcome}.
\end{split}
\label{eq:articulated-loop}
\end{equation}

The local authority used a two-link \MuJoCo{} arm, opposing fingertip
actuators, gravity, friction, a free cube, and hidden action delay.  Arm length,
motor gain, damping, starting configuration, camera position, cube size, cube
mass, and friction varied across the sealed worlds.

\subsection{Failure-driven relevance routing}

The lane forward model initially failed because gravity and articulated joint
coupling appeared as unexplained acceleration.  A private successor added
constant and position-dependent acceleration terms, but an unrestricted fit
scored only 3/4 and lost to a strong fixed-model controller at 4/4.  It was
rejected.

A selective model router then retained discovered delay when useful but fell
back to robust dynamics when nonlinear residuals were unreliable:
\begin{equation}
\begin{cases}
\text{use identified dynamics}, & \text{if residual and conditioning are credible},\\
\text{retain delay and use robust fallback}, & \text{otherwise}.
\end{cases}
\label{eq:model-router}
\end{equation}

This is the physical analogue of governed semantic routing: more learned model
content is not always better; competence includes refusing an unreliable
internal model.

\subsection{Sealed local tournament}

\begin{table}[H]
\centering
\caption{Local articulated cube-lifting tournament.}
\begin{tabular}{lccc}
\toprule
Measure & Routed \AION{} & Strong fixed ablation & Cold \\
\midrule
Successful lifts & 4/4 & 4/4 & 0/4 \\
Mean lift height & 0.08182 m & 0.08178 m & 0.00243 m \\
Privileged runtime geometry & 0 & 0 & 0 \\
Teacher actions & 0 & 0 & 0 \\
\bottomrule
\end{tabular}
\end{table}

There were 570 pre-action commitments, four of four malicious action classes
were blocked, and the simulator-neutral capsule survived restart.  \HexCore{}
promoted \texttt{procedure\_articulated\_precision\_cube\_lifting\_v1}.
The result authorized construction of an Isaac adapter and one frozen PhysX
tournament; it did not establish transfer to NVIDIA physics.

\section{Stage III: Franka Adapter and Offline Transfer Gates}

The Franka adapter was reconstructed from successful historical Isaac teacher
programs.  Forty programs were used for private construction, eight for
development criticism, and five remained sealed.  The deployed capsule was
converted to a digest-checked NumPy representation requiring no
\texttt{joblib} or \texttt{scikit-learn} installation inside the Isaac
container.

The first capsule contained an RGB cube detector, camera-to-table projection,
five cube-conditioned seven-joint waypoints, an identified action-to-motion
model, one-step inverse-dynamics corrections, and v13 fallback for admitted
out-of-distribution cases.

\begin{table}[H]
\centering
\caption{Initial offline Franka authorization gates.}
\begin{tabular}{lrr}
\toprule
Gate & Sealed result & Decision \\
\midrule
RGB cube world-position mean error & 0.000630 m & \pass \\
RGB cube world-position p90 error & 0.000907 m & \pass \\
Joint-waypoint mean error & 0.050689 rad & \pass \\
Joint-waypoint p90 error & 0.108063 rad & \pass \\
Next-motion mean error & 0.007623 rad & \pass \\
Next-motion p90 error & 0.020122 rad & \pass \\
Action-matrix condition number & 31.19 & \pass \\
Malicious contracts rejected & 6/6 & \pass \\
\bottomrule
\end{tabular}
\end{table}

These metrics authorized a paid tournament; they were never treated as
physical competence.

\section{Stage IV: First PhysX Rejection---Joint Similarity Is Not Contact}

The first frozen cloud exam used
\texttt{Isaac-Lift-Cube-Franka-v0}, seeds 99107--99118, 12 episodes per arm,
250 actions per episode, teacher actions disabled, and no runtime cube or
contact state.

\begin{table}[H]
\centering
\caption{Initial fresh PhysX tournament.}
\begin{tabular}{lrrrr}
\toprule
Arm & Strict lifts & Mean max. height & Mean return & Unsafe \\
\midrule
Joint-waypoint challenger & 0/12 & 0.055000 m & 0.6580 & 0 \\
Protected v13 & 2/12 & 0.132124 m & 15.7639 & 0 \\
\bottomrule
\end{tabular}
\end{table}

The challenger progressed through all five stored phases but never moved the
cube above its initial height.  It was rejected; v13 remained champion.  The
failure falsified the implication
\begin{equation}
\text{low joint imitation error}
\not\Rightarrow
\text{correct end-effector/object relation}
\not\Rightarrow
\text{contact or lift}.
\label{eq:falsified-joint-assumption}
\end{equation}

A seven-joint arm is redundant: many plausible joint configurations can have
small imitation error while placing or orienting the gripper incorrectly.  The
offline evaluator had rewarded resemblance to a teacher trajectory rather than
verified reduction of task-space error.

\IfFileExists{../../results/visual/precision_franka_physx_sequence.png}{
\begin{figure}[H]
\centering
\includegraphics[width=\linewidth]{../../results/visual/precision_franka_physx_sequence.png}
\caption{Representative first PhysX challenger sequence.  The robot follows a
plausible joint programme but fails to make object-moving contact.}
\end{figure}
}{}

\section{Stage V: Task-Space Servoing and Live Visual Diagnosis}

\subsection{Public Franka kinematics}

The redesign used the source-attested Franka kinematic chain to reconstruct the
hand pose from the seven bounded joint angles.  For transform $T_i(q_i)$,
\begin{equation}
T_{0h}(q)=\prod_{i=1}^{7} T_i(q_i)\,T_{7h},
\qquad
p_h(q)=\operatorname{translation}(T_{0h}(q)).
\end{equation}

Finite differences produce a positional Jacobian
\begin{equation}
J_p[:,i]\approx \frac{p_h(q+\epsilon e_i)-p_h(q)}{\epsilon}.
\end{equation}
The controller combines the RGB-derived cube position $p_c$ with a
phase-specific target height $z_k$ and computes a bounded task-space correction
\begin{equation}
\Delta q = J_p^{+}\,
\operatorname{clip}_{\delta}
\left(
\begin{bmatrix}p_{c,x}&p_{c,y}&z_k\end{bmatrix}^{\!T}-p_h(q)
\right).
\label{eq:position-servo}
\end{equation}

The identified action-to-motion model converts desired joint displacement into
a bounded actuator command.  The arm reobserves after every action.

\subsection{First live task-space result}

On the same 12-seed budget, position servoing raised mean return from 0.6580 to
0.9067 but still produced 0/12 lifts.  Protected v13 produced 1/12 on that
repetition.  Visual evidence was more informative than the scalar failure:
the arm approached the cube, arrived with unsuitable wrist/finger orientation,
disturbed the object, and then lifted empty.  Thus position was materially
improved, while contact orientation and closure timing remained unsolved.

\IfFileExists{../../results/visual/task_space_franka_physx_sequence.png}{
\begin{figure}[H]
\centering
\includegraphics[width=\linewidth]{../../results/visual/task_space_franka_physx_sequence.png}
\caption{Position-only task-space servo.  The hand reaches the object region,
but orientation and contact timing are inadequate and the cube is displaced.}
\end{figure}
}{}

\section{Stage VI: Predictive Position--Orientation Deliberation}

\subsection{Rejected private solvers}

A six-dimensional pose Jacobian was introduced:
\begin{equation}
J_{6D}(q)=
\begin{bmatrix}
J_p(q)\\J_{\omega}(q)
\end{bmatrix},
\end{equation}
where $J_{\omega}$ maps joint changes to wrist rotation.  Two private solvers
were rejected offline:
\begin{enumerate}
\item a weighted simultaneous position/orientation solver, and
\item a strict position-first null-space orientation solver.
\end{enumerate}
Both improved the combined pose score on 80\% of samples but moved closer to
the cube on only 60--64\%.  The positional safety gate was not softened.

\subsection{Prediction-filtered candidate selection}

The accepted private challenger operationalized the frame-by-frame
``time-to-think'' advantage.  At each observation it generated five wrist
correction strengths
\begin{equation}
\alpha\in\{0,0.10,0.25,0.50,1.00\},
\end{equation}
formed the corresponding desired joint corrections
\begin{equation}
\Delta q^{(\alpha)} =
\Delta q_{\text{position}}+alpha\Delta q_{\text{orientation}},
\end{equation}
and passed each through the learned actuator model to predict
$\hat q_{t+1}^{(\alpha)}$.  Any candidate expected to increase hand--cube
distance was rejected:
\begin{equation}
\left\|p^* - p_h(\hat q_{t+1}^{(\alpha)})\right\|
>
\left\|p^* - p_h(q_t)\right\|+\varepsilon
\quad\Longrightarrow\quad
\alpha\ \text{inadmissible}.
\end{equation}
Among the remaining candidates, the controller minimized
\begin{equation}
S(\alpha)=
\left\|p^* - p_h(\hat q_{t+1}^{(\alpha)})\right\|
+0.08\left\|\phi\!\left(R^*R_h(\hat q_{t+1}^{(\alpha)})^T\right)\right\|,
\label{eq:candidate-score}
\end{equation}
where $\phi(\cdot)$ is the rotation-vector map.  The policy commits only the
best admissible action; independent \PhysX{} then decides what actually occurs.

This is not a claim of access to the future.  Additional wall-clock compute is
used between discrete frames to criticize possible actions, but the next
physical observation remains the sole outcome authority:
\begin{equation}
\text{belief state}
\rightarrow \text{candidate futures}
\rightarrow \text{reject harmful predictions}
\rightarrow \text{commit one action}
\rightarrow \text{independent physics}
\rightarrow \text{prediction error}.
\label{eq:predictive-loop}
\end{equation}

\subsection{Sealed offline gate}

Five held-out successful programs yielded 25 task-space transitions.

\begin{table}[H]
\centering
\caption{Predictive pose offline gate.}
\begin{tabular}{lr}
\toprule
Measure & Result \\
\midrule
Analytic contact mean $xy$ error & 0.023822 m \\
Contact vertical-offset mean & 0.092113 m \\
Contact vertical-offset standard deviation & 0.001957 m \\
Positive position improvement & 92\% \\
Positive combined-pose improvement & 100\% \\
Mean after/before position-error ratio & 0.945199 \\
Mean predicted distance reduction per step & 0.010691 m \\
Privileged runtime inputs & 0 \\
Unsafe actions & 0 \\
Decision & \pass \\
\bottomrule
\end{tabular}
\end{table}

This gate authorized one new-seed PhysX examination only.

\section{Stage VII: First Predictive-Pose PhysX Lift}

The final frozen cohort used previously unavailable seeds 99207--99218, 12
episodes, 250 actions per episode, teacher absent, and zero privileged runtime
fields.  The policy and offline-gate digests were fixed before execution.

\begin{table}[H]
\centering
\caption{Predictive pose PhysX tournament.}
\begin{tabular}{lrrrr}
\toprule
Arm & Strict lifts & Mean max. height & Mean return & Unsafe \\
\midrule
Predictive pose challenger & 1/12 & 0.093268 m & 7.9603 & 0 \\
Protected v13 & 0/12 & 0.055000 m & 0.8913 & 0 \\
\bottomrule
\end{tabular}
\end{table}

The successful challenger episode lifted the cube to 0.514217 m.  This was the
first complete teacher-removed \PhysX{} lift produced by the new task-space
lineage.  Mean return improved by approximately $12.1\times$ relative to the
original joint-waypoint challenger.  However, the fixed promotion rule required
all of the following:
\begin{enumerate}
\item challenger strict lifts greater than 2/12;
\item challenger lifts greater than protected v13;
\item zero unsafe actions;
\item a passing offline relational gate; and
\item no teacher or privileged runtime fields.
\end{enumerate}

The challenger satisfied items 2--5 but failed item 1.  It was therefore
\reject, no procedure identifier was assigned, and
\texttt{procedure\_rgb\_phase\_bound\_episodic\_memory\_v13} remained the
protected cloud champion.  A paid cold arm was omitted after the absolute gate
failed because its result could not authorize promotion.

\IfFileExists{../../results/visual/predictive_pose_franka_physx_sequence.png}{
\begin{figure}[H]
\centering
\includegraphics[width=\linewidth]{../../results/visual/predictive_pose_franka_physx_sequence.png}
\caption{Representative predictive-pose episode-zero strip.  The sealed
tournament's successful lift occurred in episode five; only episode zero was
preselected for visual capture, preventing outcome-driven frame selection.}
\end{figure}
}{}

\section{Causal Development Ladder}

\begin{longtable}{>{\raggedright\arraybackslash}p{0.19\linewidth}
                  >{\raggedright\arraybackslash}p{0.43\linewidth}
                  >{\raggedright\arraybackslash}p{0.26\linewidth}}
\caption{What each stage established and what it falsified.}\\
\toprule
Stage & Evidence earned & Remaining or falsified assumption \\
\midrule
\endfirsthead
\toprule
Stage & Evidence earned & Remaining or falsified assumption \\
\midrule
\endhead
Precision lanes & $1{,}620.37\times$ lower error; delay identification and
disturbance recovery & Not yet articulated contact \\
Local articulated & 4/4 lift versus 0/4 cold; relevance routing & Local
\MuJoCo{} success did not establish PhysX transfer \\
Offline Franka adapter & sub-millimetre cube reconstruction and accurate joint
motion model & Joint similarity was not contact competence \\
Joint-waypoint PhysX & clean 0/12 negative against v13 2/12 & Stored joint
phases failed to preserve object-relative geometry \\
Position servo & visibly reached the cube and improved mean reward & Wrist
orientation and closure disturbed the object \\
Naive pose solvers & combined pose could improve & Orientation could not be
purchased by sacrificing approach accuracy \\
Predictive pose & five candidates criticized per frame; 92\% closer, 100\%
combined improvement offline; first live lift & Contact acquisition remained
inconsistent at 1/12 \\
\bottomrule
\end{longtable}

The principal intelligence advance is not the final score in isolation.  It is
the governed sequence
\begin{equation}
\text{observe failure}
\rightarrow \text{localize representation error}
\rightarrow \text{construct alternatives}
\rightarrow \text{reject weak variants locally}
\rightarrow \text{freeze one candidate}
\rightarrow \text{face simulator-disjoint consequences}
\rightarrow \text{retain the causal lesson}.
\end{equation}

\section{Safety, Reproducibility, and Cost Discipline}

\begin{itemize}
\item All cloud arms used the same Isaac environment, episode count, step
budget, and committed seed cohort.
\item Runtime observations contained RGB and bounded joint/gripper
proprioception only.
\item Teacher actions and privileged simulator object/contact fields were zero.
\item Every action was finite, bounded, dimension-checked, and committed before
execution.
\item Capsule and manifest digests were checked before use.
\item The protected v13 policy remained immutable throughout challenger work.
\item Development failures did not overwrite the champion.
\item Cold execution was omitted only after an absolute promotion requirement
had already failed and cold could not change the decision.
\item The NVIDIA L40S was stopped immediately after receipts and frame evidence
were copied locally.
\item No live repository or objective mutation was authorized by the physical
policy.
\end{itemize}

The efficient training strategy remains thousands of randomized local contact
episodes and only occasional frozen cloud exams.  Local randomization should
cover object pose, size, mass, friction, robot configuration, camera jitter,
action delay, grip strength, sensor noise, and forced slip.  Isaac should not
become an expensive practice loop.

\section{Next Experimental Contract: Contact Consistency}

The current bottleneck is no longer gross vision or gross arm motion.  It is
repeatable acquisition and retention of contact.  The next local curriculum
must separate the following capabilities:

\begin{enumerate}
\item \textbf{Pre-grasp alignment:} place the tool centre point above the cube
without collision.
\item \textbf{Finger straddle:} infer from visual/finger geometry whether the
cube lies between the jaws rather than merely near the wrist.
\item \textbf{Contact-timed closure:} close only when both jaws are predicted
to contact the object symmetrically.
\item \textbf{Micro-slide recovery:} after partial or one-sided contact, open
slightly, translate 1--3 cm, and reobserve.
\item \textbf{Lift verification:} raise 2--5 cm and verify that the cube moves
with the gripper.
\item \textbf{Slip response:} if visual object height or gripper/object
co-motion diverges, brake, lower, re-centre, and retry.
\item \textbf{Abandon/reapproach:} after a bounded number of failed corrections,
retreat completely and construct a new approach.
\end{enumerate}

Nominal and recovery memories must remain distinct:
\begin{equation}
\text{approach}\parallel\text{pre-grasp}\parallel\text{close}
\parallel\text{lift}\parallel\text{recover},
\end{equation}
with a router selecting one short, interruptible option at a time.  Full
eight-step recovery chunks and indiscriminate nominal/recovery pooling remain
disallowed by earlier negative results.

Before another paid exam, a frozen contact capsule should satisfy:
\begin{itemize}
\item at least 80\% held-out local success at every curriculum layer;
\item explicit improvement over the current predictive-pose capsule;
\item fresh body, camera, delay, friction, and object randomization;
\item original-fail/candidate-pass contact counterexamples;
\item recovery after forced misses and slips;
\item delayed reconstruction from the retained method rather than replayed
trajectories;
\item zero unsafe actions and explicit abstention outside admitted geometry.
\end{itemize}

Only then should a new frozen PhysX cohort compare
\begin{equation}
\text{contact-consistent challenger}
\quad\text{vs}\quad
\text{protected v13}
\quad\text{vs}\quad
\text{cold},
\end{equation}
with promotion still requiring more than 2/12 strict lifts, victory over both
controls, and no safety or authority failure.

\section{Claim Boundary}

The programme has established:
\begin{itemize}
\item transferable, uncertainty-aware motor calibration across unfamiliar
local bodies;
\item bounded articulated visual cube lifting under local MuJoCo physics;
\item a fail-closed portable Franka adapter;
\item falsification of joint-space imitation as manipulation authority;
\item object-relative position and wrist-orientation reasoning;
\item prediction-filtered candidate action selection between frames; and
\item one complete teacher-removed lift under fresh NVIDIA PhysX physics.
\end{itemize}

It has not established reliable Franka lifting, superiority over v13 across
cohorts, general manipulation, real-robot competence, Cosmos or GR00T mastery,
autonomous general apprenticeship, or AGI.  The cloud champion remains v13 and
the predictive-pose controller remains an unpromoted research challenger.

\section{Contact-Relation Learning and Recovery Arbitration}

The subsequent rapid cycle treated reliable contact as a distinct competence.
For successful teacher trajectory $i$, AION extracted the first closure event
and learned the object-relative relation
\begin{equation}
  \Delta_i = p^{\mathrm{hand}}_i - p^{\mathrm{cube}}_{i,0},
  \qquad
  \widehat{\Delta}(x,y)=W\phi(x,y)+b,
\end{equation}
where $\phi(x,y)=[1,x,y,x^2,xy,y^2]$.  The relation transferred to sealed
successful programs with mean error $0.01697$ m and p90 error $0.02955$ m.
All RGB geometry, waypoint, dynamics and relation gates passed before live
execution.

On fresh PhysX seeds $99307$--$99318$, the v2 contact controller achieved one
strict lift from twelve, tying protected v13 at one from twelve.  It therefore
failed both the absolute and comparative gates.  Trace attribution revealed a
more general defect: 56.37\% of proposals reverted to the old policy.  Contact
had moved the cube outside the narrow reset distribution in 1,171 actions, and
520 further actions were rejected because no full step was predicted to improve
alignment.  The controller was treating the consequence of its own intervention
as a reason to stop learning from that intervention.

V3 introduced table-bounded re-acquisition, re-observation after failed closure,
object-relative retry offsets, predictive step-size backtracking, and the rule
that elapsed time alone cannot authorize a phase transition.  This reduced
fallback to 6.6\%, but an excessively conservative hold rule produced 877
no-motion observations and trapped 1,372 proposals in coarse approach.  On
fresh seeds $99407$--$99418$, v3 achieved zero strict lifts while protected v13
achieved one.  V3 was rejected.

The frozen v4 proposal is the synthesis of these negative results.  It uses a
progress-conditioned coarse transition; selects the best ten-percent-scale
information action when all full steps are pessimistic; verifies attachment
with an 8 cm lift before a 25 cm lift; and converts detected slip into reopen,
re-observe and re-plan.  Its policy digest is
\texttt{738e2d02a91e57a054f08e3c4e1015ec75307968efd747ce7fc4546ed1f3994d}.
No PhysX result is claimed for v4.  The prospective $99507$--$99518$ tournament
could not start because the external Brev account had insufficient credit.
The proposal remains frozen and unpromoted.

\begin{table}[htbp]
\centering
\caption{Prospective contact-recovery generations under NVIDIA PhysX.}
\begin{tabular}{lcccc}
\toprule
Generation & AION lifts & v13 lifts & Fallback & Decision \\
\midrule
v2 contact relation & 1/12 & 1/12 & 56.37\% & reject \\
v3 re-plan/backtrack & 0/12 & 1/12 & 6.60\% & reject \\
v4 recovery synthesis & --- & --- & --- & external exam pending \\
\bottomrule
\end{tabular}
\end{table}

One execution used a container-relative evidence path and lost its artifact at
container teardown.  It was treated as invalid, assigned no score, and repeated
with the identical frozen seed cohort using an absolute mounted path.  This
distinction prevents infrastructure activity from masquerading as scientific
evidence.

\section{Artifact Manifest}

\small
\begin{longtable}{p{0.30\linewidth}p{0.64\linewidth}}
\toprule
Artifact & Repository path \\
\midrule
\endfirsthead
\toprule
Artifact & Repository path \\
\midrule
\endhead
Precision implementation & \path{backend/modules/hexcore/precision_control_apprenticeship.py} \\
Precision sealed result & \path{results/hexcore_precision_control_apprenticeship.json} \\
Articulated implementation & \path{backend/modules/hexcore/articulated_precision_cube_lifting.py} \\
Articulated sealed result & \path{results/hexcore_articulated_precision_cube_lifting.json} \\
Franka adapter builder & \path{integrations/isaac_lab/build_precision_franka_adapter.py} \\
Predictive runtime policy & \path{integrations/isaac_lab/aion_precision_franka_policy.py} \\
Task-space offline evaluator & \path{integrations/isaac_lab/evaluate_task_space_franka_gate.py} \\
Predictive CAU evaluator & \path{integrations/isaac_lab/evaluate_predictive_pose_franka_tournament.py} \\
Initial adapter result & \path{results/hexcore_precision_franka_adapter.json} \\
Initial PhysX decision & \path{results/hexcore_precision_franka_physx_tournament.json} \\
Predictive offline gate & \path{results/hexcore_task_space_franka_gate.json} \\
Predictive PhysX decision & \path{results/hexcore_predictive_pose_franka_tournament.json} \\
Contact-recovery decision & \path{results/hexcore_contact_recovery_franka_lineage.json} \\
Initial immutable receipt & \path{results/immutable/isaac_precision_franka/precision_franka_physx_receipt.json} \\
Position-servo receipt & \path{results/immutable/isaac_precision_franka/task_space_franka_physx_receipt.json} \\
Predictive-pose receipt & \path{results/immutable/isaac_precision_franka/predictive_pose_franka_physx_receipt.json} \\
Contact-relation receipt & \path{results/immutable/isaac_precision_franka/contact_consistent_franka_physx_receipt.json} \\
Re-plan receipt & \path{results/immutable/isaac_precision_franka/replan_backtracking_franka_physx_receipt.json} \\
Portable NumPy capsule & \path{results/immutable/isaac_precision_franka/precision_franka_v1.npz} \\
Initial visual strip & \path{results/visual/precision_franka_physx_sequence.png} \\
Position visual strip & \path{results/visual/task_space_franka_physx_sequence.png} \\
Predictive visual strip & \path{results/visual/predictive_pose_franka_physx_sequence.png} \\
Contact visual strip & \path{results/visual/contact_consistent_franka_physx_sequence.png} \\
Re-plan visual strip & \path{results/visual/replan_backtracking_franka_physx_sequence.png} \\
Roadmap & \path{docs/aion/AION_INTEGRATED_AGI_DEVELOPMENT_ROADMAP.md} \\
\bottomrule
\end{longtable}
\normalsize

\end{document}
